%% LyX 2.3.6.1 created this file.  For more info, see http://www.lyx.org/.
%% Do not edit unless you really know what you are doing.
\documentclass[11pt,english,11pt,a4paper,utf8]{article}
\usepackage{lmodern}
\renewcommand{\familydefault}{\rmdefault}
\usepackage[T1]{fontenc}
\usepackage{xcolor}
\usepackage{graphicx}

\makeatletter
%%%%%%%%%%%%%%%%%%%%%%%%%%%%%% User specified LaTeX commands.


%------------------------------------------------------------------------------
%	REQUIRED PACKAGES AND  CONFIGURATIONS
%------------------------------------------------------------------------------
% PACKAGES FOR TITLES
\usepackage{titlesec}
\usepackage{color}
% PACKAGES FOR LANGUAGE AND FONT
\usepackage[utf8]{inputenc}
\usepackage[english]{babel}
\usepackage[T1]{fontenc} % Font encoding

% PACKAGES FOR IMAGES
\usepackage{graphicx}
\graphicspath{{Images/}}
\usepackage{eso-pic} % For the background picture on the title page
\usepackage{subfig} % Numbered and caption subfigures using \subfloat
\usepackage{caption} % Coloured captions
\usepackage{transparent}

% STANDARD MATH PACKAGES
\usepackage{amsmath}
\usepackage{amsthm}
\usepackage{bm}
\usepackage[overload]{empheq}  % For braced-style systems of equations

% PACKAGES FOR TABLES
\usepackage{tabularx}
\usepackage{longtable} % tables that can span several pages
\usepackage{colortbl}

% PACKAGES FOR ALGORITHMS (PSEUDO-CODE)
\usepackage{algorithm}
\usepackage{algorithmic}

% PACKAGES FOR REFERENCES & BIBLIOGRAPHY
\usepackage[colorlinks=true,linkcolor=black,anchorcolor=black,citecolor=black,filecolor=black,menucolor=black,runcolor=black,urlcolor=black]{hyperref} % Adds clickable links at references
\usepackage{cleveref}
\usepackage[square, numbers, sort&compress]{natbib} % Square brackets, citing references with numbers, citations sorted by appearance in the text and compressed
\bibliographystyle{plain} % You may use a different style adapted to your field

% PACKAGES FOR THE APPENDIX
\usepackage{appendix}

% PACKAGES FOR ITEMIZE & ENUMERATES 
\usepackage{enumitem}

% OTHER PACKAGES
\usepackage{amsthm,thmtools,xcolor} % Coloured "Theorem"
\usepackage{comment} % Comment part of code
\usepackage{fancyhdr} % Fancy headers and footers
\usepackage{lipsum} % Insert dummy text
\usepackage{tcolorbox} % Create coloured boxes (e.g. the one for the key-words)

% Configuration package
\usepackage[bottom=2.0cm,top=2.0cm,left=2.0cm,right=2.0cm]{geometry}
\raggedbottom 

% Create color blue (-> manuale grafica coordinata:  https://www.polimi.it/fileadmin/user_upload/il_Politecnico/grafica-coordinata/2015_05_11_46xy_manuale_grafica_coordinata.pdf)
\definecolor{blue}{cmyk}{0.4,0.1,0,0.4}

% Custom theorem environments
\declaretheoremstyle[
  headfont=\color{blue}\normalfont\bfseries,
  bodyfont=\color{black}\normalfont\itshape,
]{colored}

\captionsetup[figure]{labelfont={color=blue}} % Set colour of the captions
\captionsetup[table]{labelfont={color=blue}} % Set colour of the captions
\captionsetup[algorithm]{labelfont={color=blue}} % Set colour of the captions

\theoremstyle{colored}
\newtheorem{theorem}{Theorem}[section]
\newtheorem{proposition}{Proposition}[section]

% Enhances the features of the standard "table" and "tabular" environments.
\newcommand\T{\rule{0pt}{2.6ex}}
\newcommand\B{\rule[-1.2ex]{0pt}{0pt}}

% Algorithm description
\newcounter{algsubstate}
\renewcommand{\thealgsubstate}{\alph{algsubstate}}
\newenvironment{algsubstates}{
    \setcounter{algsubstate}{0}%
    \renewcommand{\STATE}{%
    \stepcounter{algsubstate}%
    \Statex {\small\thealgsubstate:}\space}
    }{}
    
% Custom theorem environment
\newcolumntype{L}[1]{>{\raggedright\let\newline\\\arraybackslash\hspace{0pt}}m{#1}}
\newcolumntype{C}[1]{>{\centering\let\newline\\\arraybackslash\hspace{0pt}}m{#1}}
\newcolumntype{R}[1]{>{\raggedleft\let\newline\\\arraybackslash\hspace{0pt}}m{#1}}

% Custom itemize environment
\setlist[itemize,1]{label=$\bullet$}
\setlist[itemize,2]{label=$\circ$}
\setlist[itemize,3]{label=$-$}
\setlist{nosep}

% Create command for background pic
\newcommand\BackgroundPic{% Adding background picture
	\put(237,365){
	    \parbox[b][\paperheight]{\paperwidth}{%
	    \vfill
		\centering
		\transparent{0.4}
		\includegraphics[width=0.44\paperwidth]{raggiera_polimi}%
		\vfill}
		}
}

% Set indentation
\setlength\parindent{0pt}

% Custom title commands
\titleformat{\section}
{\color{blue}\normalfont\Large\bfseries}
{\color{blue}\thesection.}{1em}{}
\titlespacing*{\section}
{0pt}{3.3ex}{3.3ex}

\titleformat{\subsection}
{\color{blue}\normalfont\large\bfseries}
{\color{blue}\thesubsection.}{1em}{}
\titlespacing*{\subsection}
{0pt}{3.3ex}{3.3ex}

% Custom headers and footers
\pagestyle{fancy}
\fancyhf{}
      
\fancyfoot{}
\fancyfoot[C]{\thepage} % page
\renewcommand{\headrulewidth}{0mm} % headrule width
\renewcommand{\footrulewidth}{0mm} % footrule width

%%
\makeatletter
%%

\makeatother

\usepackage{babel}
\begin{document}
%\newcommand{\advisor}{Prof. Nicola Lusardi} 
% IF AND ONLY IF you need to modify the co-supervisors you also have to modify the file Configuration_files/title_page.tex (ONLY where it is marked) 
%\newcommand{\firstcoadvisor}{Name Surname} 
% insert if any otherwise comment 
%\newcommand{\secondcoadvisor}{Name Surname} 
% -> academic year 
%\newcommand{\YEAR}{2021-2022}
%
%\renewcommand{\abstract}{Insert Abstract here}
% 
\newcommand{\keywords}{Time-Tagging Module (TTM), Time-to-Digital Converter (TDC), IP-Core, Field-Programmable Gate Array (FPGA), USB 3.0}

\AddToShipoutPicture*{\BackgroundPic}

\hspace{-0.6cm}\includegraphics[width=0.6\textwidth]{Images/logo_polimi_ing_indinf}

\vspace{-1mm}

\textbf{\textcolor{blue}{\Large{}Title}}{\Large\par}

\vspace{0.5cm}

\textbf{\textsc{\textcolor{blue}{\footnotesize{}Tesi di Laurea Magistrale
in}}}{\footnotesize\par}

\textbf{\textsc{\textcolor{blue}{\footnotesize{}Electronics Engineering
- Ingegneria Elettronica}}}{\footnotesize\par}

\vspace{0.3cm}

\textbf{\large{}Chiara Muscari Tomajoli, 10560224}{\large\par}

\vspace{11pt}

\rule{1\columnwidth}{0.4pt}

\noindent\begin{minipage}[t]{1\columnwidth}%
\enskip{}%
\begin{minipage}[t]{0.24\textwidth}%
\begin{minipage}[t]{0.9\textwidth}%
\textbf{\scriptsize{}Advisor:}{\scriptsize\par}

{\scriptsize{}Prof. Nicola Lusardi}{\scriptsize\par}

\medskip{}

\textbf{\scriptsize{}Co-Advisors:}{\scriptsize\par}

{\scriptsize{}Dr. Nicola Corna}{\scriptsize\par}

{\scriptsize{}Dr. Fabio Garzetti}{\scriptsize\par}

\medskip{}

\textbf{\scriptsize{}Academic Year:}{\scriptsize\par}

{\scriptsize{}2021-2022 }{\scriptsize\par}%
\end{minipage}%
\end{minipage}\enskip{}\enskip{}%
\begin{minipage}[t]{0.74\columnwidth}%
\textbf{\textcolor{blue}{\footnotesize{}Advisor: }}\textcolor{black}{\footnotesize{}Here
goes the Abstract in English of your thesis (in article format) followed
by a list of keywords. The Abstract is a concise summary of the content
of the thesis (single page of text) and a guide to the most important
contributions included in your thesis. The Abstract is the very last
thing you write. It should be a self-contained text and should be
clear to someone who hasn\textquoteright t (yet) read the whole manuscript.
The Abstract should contain the answers to the main research questions
that have been addressed in your thesis. It needs to summarize the
motivations and the adopted approach as well as the findings of your
work and their relevance and impact. The Abstract is the part appearing
in the record of your thesis inside POLITesi, the Digital Archive
of PhD and Master Theses (Laurea Magistrale) of Politecnico di Milano.
The Abstract will be followed by a list of four to six keywords. Keywords
are a tool to help indexers and search engines to find relevant documents.
To be relevant and effective, keywords must be chosen carefully. They
should represent the content of your work and be specific to your
field or sub-field. Keywords may be a single word or two to four words. }%
\end{minipage}%
\end{minipage}

\vspace{15pt}

\begin{tcolorbox}[arc=0pt, boxrule=0pt, colback=blue!60, width=\textwidth, colupper=white]     
	\textbf{Key-words:} \keywords 
\end{tcolorbox}

\noindent\colorbox{blue}{\begin{minipage}[t]{1\textwidth}%
\textcolor{white}{\enskip{}An enhanced box with a color gradient}%
\end{minipage}}

\noindent\colorbox{cyan}{\begin{minipage}[t]{1\textwidth}%
\textcolor{white}{\enskip{}An enhanced box with a color gradient}%
\end{minipage}}

\vspace{12pt}

\include{1.Introduction/1.Introduction}

\include{2.State-of-the-Art/2.State-of-the-Art}

\include{3.Main-Work/3.Main-Work}

\include{4.Measures/4.Measures}

\include{5.Conclusions/5.Conclusions}
\end{document}
